% generated by make_tex_r1.py -- do not edit
\begin{table}[H]
\tablesize{\footnotesize}
\caption{Reflection-symmetry violation $S(t)$ (Section~\ref{sec:2d}) of the Fourier-Spectral-PINN (one run; $M=8$, $\varepsilon=0.3$) and of a vanilla PINN with the same periodic input encoding (mean~$\pm$~SD over five seeds). Data of the first version, unchanged.\label{tab:2d}}
\newcolumntype{C}{>{\centering\arraybackslash}X}
\begin{tabularx}{\textwidth}{cCC}
\toprule
\ensuremath{\boldsymbol{t}} & \textbf{Fourier-Spectral-PINN} & \textbf{Vanilla PINN (5 seeds)}\\
\midrule
0.0 & \ensuremath{8.7\times10^{-32}} & \ensuremath{(4.12\pm1.13)\times10^{-5}}\\
0.2 & \ensuremath{9.4\times10^{-32}} & \ensuremath{(4.94\pm1.16)\times10^{-5}}\\
0.5 & \ensuremath{1.0\times10^{-31}} & \ensuremath{(9.46\pm1.81)\times10^{-5}}\\
1.0 & \ensuremath{1.3\times10^{-31}} & \ensuremath{(2.22\pm0.45)\times10^{-4}}\\
\bottomrule
\end{tabularx}
\end{table}

\begin{table}[H]
\tablesize{\footnotesize}
\caption{Fourier-Spectral-PINN: reflection-symmetry violation at $t=1$ for larger mode numbers $M$ and a sharper interface ($\varepsilon=0.1$). $\sqrt{S}$ is the root-mean-square asymmetry, to be compared with the machine epsilon $2.2\times10^{-16}$.\label{tab:stress-2d}}
\newcolumntype{C}{>{\centering\arraybackslash}X}
\begin{tabularx}{\textwidth}{cCCCC}
\toprule
\ensuremath{\boldsymbol{M}} & \ensuremath{\boldsymbol{\varepsilon}} & \textbf{Epochs} & \ensuremath{\boldsymbol{S(1)}} & \ensuremath{\boldsymbol{\sqrt{S(1)}}}\\
\midrule
8 & 0.3 & 3000 & \ensuremath{1.27\times10^{-31}} & \ensuremath{3.6\times10^{-16}}\\
16 & 0.3 & 3000 & \ensuremath{1.27\times10^{-31}} & \ensuremath{3.6\times10^{-16}}\\
16 & 0.3 & 6000 & \ensuremath{1.26\times10^{-31}} & \ensuremath{3.5\times10^{-16}}\\
8 & 0.1 & 3000 & \ensuremath{3.74\times10^{-31}} & \ensuremath{6.1\times10^{-16}}\\
16 & 0.1 & 6000 & \ensuremath{3.74\times10^{-31}} & \ensuremath{6.1\times10^{-16}}\\
\bottomrule
\end{tabularx}
\end{table}
